% Part: conditional-logics
% Chapter: minimal-change-semantics
% Section: true-false

\documentclass[../../../include/open-logic-section]{subfiles}

\begin{document}

\olfileid{con}{min}{tf}

\olsection{Bebener lan Kaluputan Kontrafaktual}

Ing tuladha kang digambarake iki, kang nduweni donya anteseden
kang paling cedhak, kontrafaktual $!A \cif !B$ bener (kanthi
ora vakum) yen kabeh donya $!A$ kang paling cedhak iku donya
$!B$, kaya kang digambarake ing \olref{fig:true}.
\begin{figure}
\begin{center}
\begin{tikzpicture}[scale=.7]\small
  \spheresystem{5}
  \draw (0,0) node {$w$};
  \propositionintersect{0}{4}{40}{3.5}
  {\tikzset{proposition/.append={smooth,tension=1.4}}
  \proposition[shift={(1.6,-1)}]{90}{1}{30}{4}}
  \path (1.5,3) node[below] {$\formula{B}$};
  \path (3,0) node {$\formula{A}$};
\end{tikzpicture}
\caption{Kontrafaktual kang bener kanthi ora vakum}
\ollabel{fig:true}
\end{center}
\end{figure}
Kontrafaktual uga bener ing $w$ yen sistem bola ing saubenge~$w$
ora nduweni bola kang ngakoni $!A$ babar pisan. Ing kasus iku,
kontrafaktual bener kanthi \emph{vakum} (delengen \olref{fig:vacuous}).
\begin{figure}
\begin{center}
\begin{tikzpicture}[scale=.7]\small
  \spheresystem{5}
  \draw (0,0) node {$w$};
  \proposition{0}{6.5}{40}{3.5}
  {\tikzset{proposition/.append={smooth,tension=1.4}}
  \proposition[shift={(1.2,-1)}]{90}{1}{30}{4}}
  \path (1.5,3) node[below] {$\formula{B}$};
  \spherepos{0}{7}{node {$\formula{A}$}}
\end{tikzpicture}
\caption{Kontrafaktual kang bener kanthi vakum}
\ollabel{fig:vacuous}
\end{center}
\end{figure}

Ing kasus kanthi donya anteseden kang paling cedhak kaya iki,
kontrafaktual bisa luput kanthi rong cara. Salah sijine yaiku
yen ora kabeh donya $!A$ kang paling cedhak iku donya $!B$,
nanging sawetara iya. Ing kasus iki, $!A \cif \lnot !B$ uga
luput (delengen \olref{fig:false}).
\begin{figure}
\begin{center}
\begin{tikzpicture}[scale=.7]\small
  \spheresystem{5}
  \draw (0,0) node {$w$};
  \propositionintersect{0}{4}{40}{3.5}
  {\tikzset{proposition/.append={smooth,tension=1.4}}
  \proposition[shift={(1.6,0)}]{90}{1}{30}{3}}
  \path (1.5,3) node[below] {$\formula{B}$};
  \path (3,0) node {$\formula{A}$};
\end{tikzpicture}
\caption{Kontrafaktual luput, kondisional kanthi konsekuen kang dinegasi uga luput}
\ollabel{fig:false}
\end{center}
\end{figure}
Yen donya-donya $!A$ kang paling cedhak ora duwe irisan karo
donya-donya $!B$ babar pisan, mula $!A \cif !B$ luput.
Nanging ing kasus iki, kabeh donya $!A$ kang paling cedhak
iku donya $\lnot !B$, mula $!A \cif \lnot !B$ bener
(delengen \olref{fig:false-opposite}).
\begin{figure}
\begin{center}
\begin{tikzpicture}[scale=.7]\small
  \spheresystem{5}
  \draw (0,0) node {$w$};
  \propositionintersect{0}{4}{40}{3.5}
  {\tikzset{proposition/.append={smooth,tension=1.4}}
  \proposition[shift={(-1,0)}]{90}{1}{30}{3}}
  \path (-1,3) node[below] {$\formula{B}$};
  \path (3,0) node {$\formula{A}$};
\end{tikzpicture}
\caption{Kontrafaktual luput, kondisional kanthi konsekuen kang dinegasi bener}
\ollabel{fig:false-opposite}
\end{center}
\end{figure}

Beda karo kondisional ketat kang wis dirembug, kontrafaktual bisa
kontingen. Gatekna model bola ing \olref{fig:contingent}.
Kabeh donya $!A$ kang paling cedhak karo~$u$ iku donya $!B$,
mula $\mSat{M}{!A \cif !B}[u]$. Nanging ana donya $!A$
kang paling cedhak karo~$v$ lan dudu donya $!B$, mula
$\mSat/{M}{!A \cif !B}[v]$.
\begin{figure}
\begin{center}
\begin{tikzpicture}[scale=.6]\small
  \spheresystem{7}
  \spheresystem[shift={(0:3.1)},dashed]{4}
  \propositionintersect{45}{4}{40}{4.5}
  \begin{scope}
    \clip \propositionplot{45}{4}{40}{4.5} ;
    \spherelayer[shift={(0:3.1)}]{3}
  \end{scope}
  \proposition[shift={(1.4,-.5)}]{90}{1}{35}{5}
  \path (0,0) node {$u$};
  \path (0:3.1) node {$v$};
  \spherepos{35}{9}{node {$\formula{A}$}}
  \spherepos{80}{9}{node {$\formula{B}$}}
\end{tikzpicture}
\end{center}
\caption{Kontrafaktual kontingen}
\ollabel{fig:contingent}
\end{figure}
\end{document}
